% Abhakseite „Das kann ich“ – Zone nur im Gesamt (\mitzone)
\begin{abhakseite}
\ifdefined\mitzone
\abhakgruppe{Kennst du schon}
\abhak{1}{Ich kann negative Zahlen addieren und subtrahieren.}
\abhak{2}{Ich finde den Fehler beim Rechnen mit negativen Zahlen.}
\abhak{3}{Ich kann Zahlen mit Vorzeichen malnehmen.}
\abhak{4}{Ich kann mit Dezimalzahlen und einfachen Brüchen rechnen.}
\abhak{5}{Ich kann Punkt vor Strich rechnen.}
\fi
\abhakgruppe{Einheit 1 · Terme aufstellen und berechnen}
\abhak{6–7}{Ich kann einen Term aufstellen.}
\abhak{8}{Ich finde auch in Aufgaben aus der Prüfung den passenden Term.}
\abhak{9}{Ich kann Termwerte berechnen.}
\abhak{10}{Ich kann zu einem Term eine passende Situation erfinden.}
\abhak{11}{Ich finde den Fehler beim Aufstellen eines Terms.}
\abhak{12}{Ich kann begründen, ob ein Term zum Text passt.}
\abhak{13}{Ich kann Terme und Termbäume einander zuordnen.}
\abhak{14}{Ich kann Terme für Aufgaben aus dem Alltag aufstellen.}
\abhakgruppe{Einheit 2 · Terme zusammenfassen}
\abhak{15}{Ich erkenne gleichartige Glieder.}
\abhak{16}{Ich kann die Vorzahl eines Gliedes ablesen.}
\abhak{17}{Ich kann die Glieder eines Terms mit ihrem Vorzeichen aufschreiben.}
\abhak{18–19}{Ich kann Terme zusammenfassen.}
\abhak{20}{Ich kann auch Aufgaben aus der Prüfung zum Zusammenfassen lösen.}
\abhak{21}{Ich kann Terme malnehmen.}
\abhak{22}{Ich kann auch Produkte mit Vorzeichen und zwei Variablen berechnen.}
\abhak{23}{Ich finde den Fehler beim Zusammenfassen.}
\abhak{24}{Ich kann begründen, ob man Glieder zusammenfassen kann.}
\abhak{25}{Ich kann Terme zu Figuren aufstellen und zusammenfassen.}
\abhak{26}{Ich kann Terme im Alltag aufstellen und zusammenfassen.}
\abhakgruppe{Einheit 3 · Klammern auflösen}
\abhak{27}{Ich erkenne, was vor einer Klammer steht.}
\abhak{28–29}{Ich kann Klammern auflösen.}
\abhak{30}{Ich kann auch zwei Klammern in einem Term auflösen und zusammenfassen.}
\abhak{31}{Ich finde den Fehler beim Auflösen von Klammern.}
\abhak{32}{Ich kann begründen, wie man eine Klammer richtig auflöst.}
\abhak{33}{Ich kann Klammern mit einer Tabelle oder einer Figur auflösen.}
\abhak{34}{Ich kann Terme mit Klammer im Alltag aufstellen und auflösen.}
\abhakgruppe{Einheit 4 · Ausklammern}
\abhak{35}{Ich erkenne, welcher Faktor in jedem Glied steckt.}
\abhak{36}{Ich kann jedes Glied als Produkt mit einem Faktor schreiben.}
\abhak{37}{Ich kann einen vorgegebenen Faktor ausklammern.}
\abhak{38}{Ich kann ausklammern.}
\abhak{39}{Ich kann auch Terme mit zwei Variablen so weit wie möglich ausklammern.}
\abhak{40}{Ich finde den Fehler beim Ausklammern.}
\abhak{41}{Ich kann begründen, ob richtig ausgeklammert wurde.}
\abhak{42}{Ich kann einen Term mit Klammer in Worten oder als Figur beschreiben.}
\abhak{43}{Ich kann mit Ausklammern Aufgaben aus dem Alltag lösen.}
\end{abhakseite}
